\begin{obereflection}
\textbf{Latihan 16.1}
1. Sebutkan empat layanan kriptografi yang disebut modul pengantar.
2. Jelaskan makna nir-penyangkalan. Berikan contoh yang disebut modul.
3. Pada tanda tangan digital, kunci mana yang dipakai untuk menghasilkan tanda
   tangan? Kunci mana yang dipakai untuk memverifikasi?
4. Sebutkan algoritma tanda tangan digital yang disebut dalam modul.
5. Standar apa yang disebut modul RSA sebagai pengatur penggunaan RSA? Apa saja
   isinya?
\end{obereflection}
